\documentclass[lualatex, paper={90mm, 160mm}]{jlreq}
\usepackage{tikz}
% luatexja は jlreq[lualatex] が自動的に読み込む
% lltjext: 縦書きminipage (<t>オプション) のために必要
\usepackage{lltjext}
\usetikzlibrary{calc}

% 和色の定義
\definecolor{kinari}{HTML}{F6E5CC} % 生成色（背景用）
\definecolor{aiiro}{HTML}{165E83}  % 藍色（模様・枠線用）

\begin{document}
\pagestyle{empty}

\begin{tikzpicture}[remember picture, overlay]

  % ━━━━━━━━━━━━━━━━━━━━━━━━━━━━━━━━━━
  % 1. 背景色
  % ━━━━━━━━━━━━━━━━━━━━━━━━━━━━━━━━━━
  \fill[kinari!40] (current page.south west) rectangle (current page.north east);

  % ━━━━━━━━━━━━━━━━━━━━━━━━━━━━━━━━━━
  % 2. 背景模様: 市松模様（上下）
  % ━━━━━━━━━━━━━━━━━━━━━━━━━━━━━━━━━━
  \begin{scope}[opacity=0.15, fill=aiiro]
    % 上部の市松模様
    \foreach \x in {0, 2, ..., 18} {
      \foreach \y in {0, 2, 4} {
        \fill ($(current page.north west) + (\x*5mm, -\y*5mm)$) rectangle ++(5mm, -5mm);
        \fill ($(current page.north west) + (\x*5mm+5mm, -\y*5mm-5mm)$) rectangle ++(5mm, -5mm);
      }
    }
    % 下部の市松模様
    \foreach \x in {0, 2, ..., 18} {
      \foreach \y in {0, 2, 4} {
        \fill ($(current page.south west) + (\x*5mm, \y*5mm)$) rectangle ++(5mm, 5mm);
        \fill ($(current page.south west) + (\x*5mm+5mm, \y*5mm+5mm)$) rectangle ++(5mm, 5mm);
      }
    }
  \end{scope}

  % ━━━━━━━━━━━━━━━━━━━━━━━━━━━━━━━━━━
  % 3. 二重枠線（和色に調整）
  % ━━━━━━━━━━━━━━━━━━━━━━━━━━━━━━━━━━
  \draw[line width=1.2pt, color=aiiro!90]
    ($(current page.south west) + (6mm, 6mm)$)
    rectangle ($(current page.north east) - (6mm, 6mm)$);
  \draw[line width=0.3pt, color=aiiro!90]
    ($(current page.south west) + (7.5mm, 7.5mm)$)
    rectangle ($(current page.north east) - (7.5mm, 7.5mm)$);

  % ━━━━━━━━━━━━━━━━━━━━━━━━━━━━━━━━━━
  % 4. タイトル（縦書き）
  % ━━━━━━━━━━━━━━━━━━━━━━━━━━━━━━━━━━
  \node[anchor=north] at ($(current page.center) + (15mm, 50mm)$) {
    \begin{minipage}<t>{15em}
      \begin{center}
        {\fontsize{28pt}{40pt}\selectfont \color{black!90}走れメロス\par}
      \end{center}
    \end{minipage}
  };

  % ━━━━━━━━━━━━━━━━━━━━━━━━━━━━━━━━━━
  % 5. 作者名（縦書き）
  % ━━━━━━━━━━━━━━━━━━━━━━━━━━━━━━━━━━
  \node[anchor=north] at ($(current page.center) + (-15mm, -10mm)$) {
    \begin{minipage}<t>{10em}
      \begin{center}
        {\fontsize{16pt}{24pt}\selectfont \color{black!90}太宰 治\par}
      \end{center}
    \end{minipage}
  };

\end{tikzpicture}
\end{document}